\chapter{Wave Equation in Solids}

\section*{Introduction}

When an earthquake strikes, seismic waves travel through the Earth's crust at thousands of meters per second. P-waves (compressional) arrive first, then S-waves (shear), then surface waves. By measuring the arrival times at different seismograph stations, scientists locate the earthquake's origin and map the Earth's interior structure---all using the wave equation in solids.

\section*{Fundamental Equation Used}

\[
\rho \frac{\partial^2 u_i}{\partial t^2} = C_{ijkl} \frac{\partial^2 u_k}{\partial x_l}
\]
\section*{Assumptions}

\textbf{Assumptions:} The solid is elastic (stress proportional to strain). Small deformations (linear elasticity). The elastic constants $C_{ijkl}$ are uniform. No damping or anisotropic absorption.


\textbf{Limitations:} Does not account for plastic deformation (permanent set after large strains). Breaks down at very high frequencies (phonon effects). Does not include thermal expansion coupling or piezoelectric effects.


\section*{Mathematical Derivation}

\textbf{Step 1: The Physical Setup.}

Consider a small volume element of an elastic solid with density $\rho$. When the solid is deformed, each material point $\mathbf{r}$ is displaced to $\mathbf{r} + \mathbf{u}(\mathbf{r}, t)$, where $\mathbf{u}$ is the displacement field. We apply Newton's second law to this volume element.

\textbf{Step 2: Stress and the Equation of Motion.}

The force on a volume element $dV$ due to internal stresses is:

\begin{equation}
F_i = \int_{\partial V} \sigma_{ij}\,n_j\,dA,
\end{equation}

where $\sigma_{ij}$ is the stress tensor and $n_j$ is the outward unit normal. By the divergence theorem:

\begin{equation}
F_i = \int_V \frac{\partial \sigma_{ij}}{\partial x_j}\,dV.
\end{equation}

For a small element, the equation of motion (Newton's second law) is:

\begin{equation}
\rho\,\frac{\partial^2 u_i}{\partial t^2} = \frac{\partial \sigma_{ij}}{\partial x_j} + f_i^{\text{body}},
\end{equation}

where $f_i^{\text{body}}$ is the body force per unit volume. Neglecting body forces:

\begin{equation}
\rho\,\frac{\partial^2 u_i}{\partial t^2} = \frac{\partial \sigma_{ij}}{\partial x_j}. \tag{1}
\end{equation}

\textbf{Step 3: The Elastic Constitutive Law.}

For a linear elastic solid, the stress is proportional to the strain (generalized Hooke's law):

\begin{equation}
\sigma_{ij} = C_{ijkl}\,\varepsilon_{kl},
\end{equation}

where $C_{ijkl}$ is the fourth-order elasticity tensor and $\varepsilon_{kl}$ is the strain tensor. For a general anisotropic solid, $C_{ijkl}$ has up to 21 independent components (for triclinic symmetry).

\textbf{Step 4: Strain in Terms of Displacement.}

The infinitesimal strain tensor is related to the displacement gradient by:

\begin{equation}
\varepsilon_{kl} = \frac{1}{2}\left(\frac{\partial u_k}{\partial x_l} + \frac{\partial u_l}{\partial x_k}\right).
\end{equation}

Substituting into the constitutive law:

\begin{equation}
\sigma_{ij} = C_{ijkl}\,\frac{1}{2}\left(\frac{\partial u_k}{\partial x_l} + \frac{\partial u_l}{\partial x_k}\right) = C_{ijkl}\,\frac{\partial u_k}{\partial x_l},
\end{equation}

where we used the symmetry $C_{ijkl} = C_{ijlk}$ to replace $\frac{1}{2}(\partial_l u_k + \partial_k u_l)$ by $\partial_l u_k$.

\textbf{Step 5: The Elastic Wave Equation.}

Substitute into the equation of motion (1):

\begin{equation}
\rho\,\frac{\partial^2 u_i}{\partial t^2} = C_{ijkl}\,\frac{\partial^2 u_k}{\partial x_l\,\partial x_j} = C_{ijkl}\,\frac{\partial^2 u_k}{\partial x_j\,\partial x_l}.
\end{equation}

This is the \textbf{elastic wave equation in a solid}:

\begin{equation}
\boxed{\rho\,\frac{\partial^2 u_i}{\partial t^2} = C_{ijkl}\,\frac{\partial^2 u_k}{\partial x_l\,\partial x_j}.}
\end{equation}

\textbf{Step 6: Plane Wave Solutions and Dispersion.}

Assume a plane wave solution $u_k = A_k\,e^{i(k_j x_j - \omega t)}$. Substituting:

\begin{equation}
-\rho\,\omega^2\,A_i = -C_{ijkl}\,k_j\,k_l\,A_k.
\end{equation}

Rearranging:

\begin{equation}
\left(C_{ijkl}\,k_j\,k_l - \rho\,\omega^2\,\delta_{ik}\right)A_k = 0.
\end{equation}

For non-trivial solutions, the determinant must vanish:

\begin{equation}
\det\left(C_{ijkl}\,k_j\,k_l - \rho\,\omega^2\,\delta_{ik}\right) = 0.
\end{equation}

This is a $3 \times 3$ eigenvalue problem that yields three wave speeds---corresponding to one quasi-longitudinal (P-wave) and two quasi-transverse (S-wave) modes for each propagation direction.

\textbf{Step 7: Isotropic Solid.}

For an isotropic solid, $C_{ijkl} = \lambda\,\delta_{ij}\delta_{kl} + \mu(\delta_{ik}\delta_{jl} + \delta_{il}\delta_{jk})$, and the wave equation simplifies to:

\begin{equation}
\rho\,\frac{\partial^2 u_i}{\partial t^2} = (\lambda + \mu)\frac{\partial^2 u_k}{\partial x_i\,\partial x_k} + \mu\,\frac{\partial^2 u_i}{\partial x_k^2}.
\end{equation}

Taking the divergence and curl of this equation decouples the modes:

\begin{itemize}
\item P-waves (compressional): $v_P = \sqrt{(\lambda + 2\mu)/\rho}$,
\item S-waves (shear): $v_S = \sqrt{\mu/\rho}$.
\end{itemize}

Since $\lambda + 2\mu > \mu > 0$, P-waves always propagate faster than S-waves, explaining the order of seismic wave arrivals.

\section*{Characteristics of the Equation}

Two types of bulk waves: P-waves (compressional, $v_P = \sqrt{(\lambda + 2\mu)/\rho}$) and S-waves (shear, $v_S = \sqrt{\mu/\rho}$). P-waves are always faster than S-waves. At surfaces, Rayleigh waves and Love waves propagate with different speeds and polarization.
\section*{Solution Methods}

Separation of variables in Cartesian, cylindrical, and spherical coordinates. Normal mode analysis (plates, beams, shells). Finite element methods for complex geometries. Spectral methods for dispersion relations.
\section*{Verifications}

Ultrasonic pulse-echo measurements give wave speeds, from which elastic constants are calculated. Resonance ultrasound spectroscopy measures natural frequencies to determine elastic moduli. Seismic data is compared against theoretical dispersion curves.
\section*{Applications}

Seismology (earthquake location and Earth structure), ultrasonic non-destructive testing (detecting cracks in bridges and aircraft), acoustic microscopy, sonar, geological exploration (oil and gas prospecting), architectural acoustics.
\section*{What Richard Feynman Would Have Asked}

\subsection*{1. The Plain-English Translation}
\textit{``What is this equation telling me about how sound moves through a solid?''}

The Core Meaning: When you tap one end of a steel rod, the atoms at that end push against their neighbors, which push against their neighbors, and a compression wave travels through the material. The wave equation in solids quantifies this: the speed depends on how stiff the atomic bonds are (elastic constants) and how heavy the atoms are (density). Stiffer bonds and lighter atoms mean faster waves.

The Metaphor: Think of a solid as a 3D network of springs connecting atoms. Tapping one end compresses the springs there, and the compression propagates through the network. Stiff springs (strong bonds) transmit the push quickly; light atoms (low density) accelerate faster, so the wave propagates faster.

\subsection*{2. The Localized Mechanics}
\textit{``At the atomic level, what creates the wave?''}

He would press you on the connection between atomic bonds and elastic constants. The elastic constants $C_{ijkl}$ are determined by the second derivatives of the interatomic potential at equilibrium spacing. Stretching the lattice (positive strain) increases the potential energy; the elastic constants measure how much. The wave equation is Newton's second law applied to each atom, coupled to its neighbors through the interatomic forces.

The Smoothness Check: He would ask why the continuum approximation works at all. The interatomic spacing is $\sim 10^{-10}$ m, while typical wavelengths are $\sim 10^{-3}$ m or larger. When the wavelength is much larger than the atomic spacing, the discrete lattice looks like a smooth continuum, and the wave equation holds.

\subsection*{3. Testing the Extreme Limits}
\textit{``What happens at extreme frequencies or wavelengths?''}

The Phonon Limit: When the wavelength approaches the interatomic spacing ($\lambda \sim a$), the continuum approximation breaks down. The wave speed becomes frequency-dependent (dispersion), and the concept of a continuous medium fails. The vibrations are then described as quantized phonons.

The Shock Wave Limit: At very large amplitudes, the linear elasticity assumption fails. Compression waves steepen into shock waves (discontinuities in stress and density). The nonlinear wave equation is needed, and energy is dissipated at the shock front.

The Earth's Core: Seismic waves slow down and refract as they pass through the Earth's mantle and core, revealing the planet's layered structure. S-waves disappear at the liquid outer core boundary, confirming its liquid state.

\subsection*{4. Dimensionality and Geometry}
\textit{``How does the geometry of the solid affect the wave speed?''}

He would point out that in an infinite solid (bulk waves), geometry doesn't matter---the wave speed is determined by material properties alone. But in thin plates, rods, and shells, boundary conditions create guided waves (Lamb waves, Love waves) with dispersion---the wave speed depends on frequency and geometry. The same material can have different wave speeds depending on its shape.

\subsection*{5. Cross-Disciplinary Connection}
\textit{``Where else does this same elastic wave equation appear?''}

The elastic wave equation appears in: lattice dynamics (phonons in crystals),seismology (seismic exploration for oil), medical ultrasound (bone density measurements), non-destructive testing (detecting flaws in pipelines), and even in the early universe (primordial gravitational waves are elastic waves in spacetime itself).
\section*{FAQ}

\textbf{Q: Why can't S-waves travel through liquids?}\\
\textbf{A:} Shear waves require shear stiffness ($\mu > 0$). Liquids have $\mu = 0$ (no resistance to shear deformation), so S-waves cannot propagate. This is how seismologists know the Earth's outer core is liquid: S-waves are blocked.

\textbf{Q: Why does sound travel faster in steel than in air?}\\
\textbf{A:} Steel is much stiffer ($E \approx 200$ GPa) than air ($B \approx 10^5$ Pa), despite being much denser. The stiffness-to-density ratio is much higher in solids, so $v = \sqrt{E/\rho}$ is much larger.
\section*{Important Bibliography}

\begin{enumerate}
\item Achenbach, J.D., \textit{Wave Propagation in Elastic Solids} (Elsevier, 1973), Chapters 1--4.
\item Graff, K.F., \textit{Wave Motion in Elastic Solids} (Dover, 1991), Chapters 3--5.
\item Shearer, P., \textit{Introduction to Seismology}, 2nd ed. (Cambridge University Press, 2009), Chapters 3--4.
\item Ewing, M., Jardetzky, W., and Press, F., \textit{Elastic Waves in Layered Media} (McGraw-Hill, 1957), Chapters 2--3.
\end{enumerate}


